\begin{figure*}[t]
\centering
\resizebox{\textwidth}{!}{%
\begin{tikzpicture}[>=Stealth, font=\small]

% ---- Styles ----
\tikzset{
  stepbox/.style={
    draw=black!30, rounded corners=4pt,
    minimum width=3.8cm, minimum height=2.2cm,
    align=center, line width=0.5pt,
    inner sep=5pt, fill=white
  },
  iobox/.style={
    draw=black!25, rounded corners=3pt,
    minimum height=0.8cm, align=center,
    fill=white, line width=0.5pt, inner sep=3pt
  },
  arr/.style={->, thick, black!35},
}

% ============ INPUT ============
\node[iobox, minimum width=2.3cm] (seed) at (-3.8, 0.55) {
  Seed Rubric $\rubric_0$
};
\node[iobox, minimum width=2.7cm] (data) at (-3.8, -0.55) {
  Training Data $\{(x_i, y_i)\}$
};

% ============ STEP I: SCORE ============
\node[stepbox] (score) at (0, 0) {
  \textbf{I.~Score}\\[3pt]
  Score each essay $x_i$\\
  with current rubric $\rubric$\\[2pt]
  {\footnotesize $(\hat{y}_i,\; z_i) = \mathrm{LLM}(\rubric,\; x_i)$}
};

% ============ STEP II: REFLECT ============
\node[stepbox] (reflect) at (5.5, 0) {
  \textbf{II.~Reflect}\\[3pt]
  Collect scoring errors\\[1pt]
  {\footnotesize $\mathcal{E} = \bigl\{(x_i, y_i, \hat{y}_i, z_i) \mid \hat{y}_i \neq y_i\bigr\}$}\\[1pt]
  {\scriptsize Balanced sampling by score level}
};

% ============ STEP III: REVISE ============
\node[stepbox] (revise) at (12, 0) {
  \textbf{III.~Revise}\\[3pt]
  Revise rubric based on\\
  errors \textbf{+ rationales}\\[2pt]
  {\footnotesize $\rubric' = \procname{ReviseRubric}(\rubric,\; \tilde{\mathcal{E}})$}
};

% ============ ERROR SET ILLUSTRATION ============
\node[draw=black!25, rounded corners=4pt, fill=white,
  inner sep=6pt] (errorset) at (8.75, -3.3) {
  \begin{minipage}{4.6cm}
    \centering\textbf{Error Set} $\mathcal{E}(\rubric)$\\[4pt]
    \fcolorbox{black!10}{black!2}{\parbox{4.1cm}{\scriptsize
      \textbf{Essay:} Paul Merson has restarted his row with\ldots\\
      \textbf{Rationale:} The essay provides minimal elaboration\ldots\\
      \textcolor{red!70!black}{\textbf{Predicted: 3}} \enspace \textbf{Human: 5}%
    }}\\[3pt]
    \fcolorbox{black!10}{black!2}{\parbox{4.1cm}{\scriptsize
      \textbf{Essay:} Dear local newspaper, I think effects of\ldots\\
      \textbf{Rationale:} The response takes a clear position\ldots\\
      \textcolor{red!70!black}{\textbf{Predicted: 5}} \enspace \textbf{Human: 4}%
    }}\\[3pt]
    $\cdots\cdots$
  \end{minipage}
};

% ============ TOP-K SELECTION ============
\node[iobox, minimum width=3.2cm] (topk) at (5.5, -6) {
  Top-$K$ selection by QWK
};

% ============ OUTPUT ============
\node[iobox, minimum width=2.8cm] (best) at (9.5, -6) {
  \textbf{Best Rubric} $\rubric_{\mathrm{best}}$
};

% ============ ARROWS ============
% Input -> Score
\draw[arr] (seed.east) -- (score.west |- seed);
\draw[arr] (data.east) -- (score.west |- data);

% Score -> Reflect
\draw[arr] (score) -- (reflect);

% Reflect -> Error Set (diagonal)
\draw[arr] (reflect.south) -- (errorset.north west);

% Error Set -> Revise (diagonal)
\draw[arr] (errorset.north east) -- (revise.south);

% Revise -> Top-K (route right of Error Set)
\draw[arr] (revise.south east) |- (topk.east);

% Top-K -> Output
\draw[arr] (topk) -- (best);

% Loop: Top-K -> Score (iteration)
\draw[arr, dashed] (topk.west) -| (score.south);

% Iteration label
\node[font=\small, text=black!45] at (1.5, -6.4) {$\times\, T$ iterations};

\end{tikzpicture}%
}
\caption{Overview of our Reflect-and-Revise framework. Starting from a seed rubric and training essays with human scores, the method iteratively: (\textbf{I})~scores each essay with the current rubric, producing predictions and chain-of-thought rationales; (\textbf{II})~reflects on scoring errors by collecting cases where model predictions disagree with human labels; and (\textbf{III})~revises the rubric based on sampled error cases and their rationales. A top-$K$ selection step retains the best-performing candidates by QWK for the next iteration.}
\label{fig:overview}
\end{figure*}
